%!TEX root = notes.tex
% Notes from --
% Typed by: Durga

\section{Introduction}

In this section we are going to discuss some of the prerequisite necessary for this course.

\noindent\textbf{Homology Group:}
Homology is a covariant functor from category of topological space to the category of graded abelian groups
\[H_\ast: \topsp \to \cx_{\geq 0} \to \GrAb_{\geq 0}.\]
For a space \(X\) we get a graded abelian group \(\displaystyle H_\ast(X) = \bigoplus_{n\geq0} H_n(X)\).
The group \(H_n(X)\) is called the \(n\)-th Homology group.
The group \(H_0(X)\) is the free abelian group generated by the number of path components of \(X\).
For more details refer to \cite{Hatcher2002}.

\noindent\textbf{Cohomology Ring:}
The cohomology is a contravarient functor  from category of topological space to the category of graded commutative rings \(\GrRing_{\geq 0}\),
\[H^\ast: \topsp^{\op} \to \GrRing_{\geq 0}.\]
For a space \(X\) we get a graded ring \(\displaystyle H^\ast(X) = \bigoplus_{n\geq0} H^n(X)\).
This ring has a multiplication defined as the cup product in \cite{Hatcher2002}, this multiplication is graded commutative. 
For more details refer to \cite{Hatcher2002}.

\noindent\textbf{Homotopy Group:}
For a pointed topological space \((X,x_0)\) we get a collection of groups \(\{\pi_n(X,x_0)\}_{n\geq 0}\).
The group \(pi_n(X,x_0)\) is called \(n\)-th homotopy group, \(n\geq 1\).
The \(\pi_0(X)\) is the set of path connected component of \(X\).
The group \(\pi_n(X)\) is abelian for \(n\geq 2\).
This forms a functor 
\[\pi_\ast: \topsp \to \Set \times \Grp \times \prod_{n\geq 2} \Ab.\]
If \(X\) is path-connected then, \(\pi_n(X,x_0)\) are independent of the base point \(x_0\).
For more details refer to \cite{Hatcher2002}.

All of the above three functors are homotopy invariant, but converse is not true, i.e., if \(X\) and \(Y\) has same homology, homotopy or cohomology groups then it does not implies that \(X\) and \(Y\) are homotopic.
In rational homotopy theory we get a converse but with some restriction: 
We are going to see a functor \(A^\ast\) from the category of simplicial complexes to category of \emph{Commutative Differential Graded \(\mathbb Q\)-Algebra} (\(\CDGA\)). 
Given a space \(X\) we are going to get a commutative differential graded \(\mathbb Q\)-algebra \(A^\ast(X)\) whose cohomology is \(H^\ast(X,\mathbb Q)\).
If there is a quasi isomorphism  between \(A^\ast(X)\) and \(A^\ast(Y)\) (i.e., it preserves cohomology groups) then \(X\) are \(Y\) are "\emph{rationally homotopy equivalent}", we denote it as "\(X \simeq_\mathbb Q Y\)".
This result is due to \emph{Sullivan, Dennis} refer to \cite{PMIHES_1977__47__269_0}

\begin{definition}[Rational Homotopy Equivalence]
	For \(X\) and \(Y\) simply connected topological space, a continuous function \(f: X \to Y\) is called a \emph{rational homotopy equivalence} if it satisfy the following equivalent conditions: 
	\begin{enumerate}
		\item the induced map \(f_\ast \otimes \mathbb Q: \pi_\ast(X) \otimes \mathbb Q \to \pi_\ast(X) \otimes \mathbb Q\) is an isomorphism.
		\item The induced map is an isomorphism on homology group with rational coefficients 
		\[f_\ast: H_\ast(X,\mathbb Q) \xrightarrow{\cong} \ast(Y,\mathbb Q),\]
		\item The induced map is an isomorphism on rational cohomology groups
		\[f^\ast : H^\ast(Y, \mathbb Q) \xrightarrow{\cong} H^\ast(X, \mathbb Q).\]
		\item it induces a weak homotopy equivalence on the "rationalization".
	\end{enumerate}
	We are going to know more about rationalization later.
\end{definition}

A natural question arises that is "how much weaker or stronger the rationally homotopy equivalent is compared to the homotopy equivalence?", and here are some examples to understand that.

\begin{example}
	The space of real projective space \(\rproj{2}\) is not homotopy equivalent to sphere \(\sphere{2}\), because they don't have isomorphic homotopy groups. 
	Since the covering map \(\pi: \sphere{2} \to \rproj{2}\) is a rational homotopy equivalence, they both are rationally equivalent.
\end{example} 

\begin{definition}[Eilenberg–MacLane space]
	For an abelian group \(G\) and \(n \geq 2\) there is a space \(\emspace{G}{n}\) called Eilenberg–MacLane space such that 
	\[\pi_i(\emspace{G}{n}) = \begin{cases}
		G, & i = n\\
		0, & \text{otherwise.}
	\end{cases}\]
\end{definition}

\begin{example}
	Two space \(\sphere{2}\) and \(\sphere{3} \times \emspace{G}{n}\) for any finite abelian group \(G\) and \(n \geq 0\) has isomorphic homotopy groups but their homology groups \(H_n(-)\) with rational coefficients are not isomorphic for \(n = 3, 4, 6, 8 \dots\), hence they are not rational homotopy equivalent.
\end{example}

\begin{example}
	The space \(\emspace{\mathbb Z/2}{2}\) and one point space have trivial \(\pi_i(-) \otimes \mathbb Q\) for all \(i \geq 0\), but they are not homotopy equivalent because they have non-isomorphic \(\pi_1(-)\).
\end{example}

For real smooth manifold \(M\), we can use de-Rham cochain complex \(C_{DR}^\ast(M)\) to compute its cohomology 
\[H_{DR}^\ast(M) \cong H^\ast(M).\]
In this case a quasi isomorphism between \(C^\ast_{DR}(M)\) and \(C^\ast_{DR}(N)\) implies \(M\) real homotopic to \(N\).
 
The overview of the process that we are going to follow is
\begin{center}
	\[\begin{tikzcd}
		{\{\text{Simplical Complex}\}} && {\{\CDGA\}} \\
		\\
		{\{\mathbb Q -\text{Space}\}} \\
		\\
		{\{\text{Postnikov Tower}/Q\}} && {\{\text{Minimal } \CDGA\}}
		\arrow["{\text{PL forms}}", from=1-1, to=1-3]
		\arrow["{\text{Localization at } 0}"', from=1-1, to=3-1]
		\arrow["{\text{Minimal Model}}", from=1-3, to=5-3]
		\arrow["{\text{Postnikov Tower}}"', from=3-1, to=5-1]
		\arrow[from=5-1, to=5-3]
		\arrow[from=5-3, to=5-1]
	\end{tikzcd}\]
\end{center}

Over the integers, the cochain complex \(C^\ast(X,\mathbb Z)\) of a space \(X\) forms a Differential Graded Algebra \((\mathsf(DGA))\) that is graded commutative up to homotopy.
It is not possible to replace these by a Commutative Differential Graded Algebra \(\mathsf{CDGA}\).
The nontrivial "Steenrod Operations" obstructs us to do so.
To encode the full homotopy type of \(X\) one needs some extra data of \(E_\infty\)-algebra structure on \(C^\ast(X,\mathbb Z)\), \cite{PMIHES_2006__103__213_0}.


